\section{An alternative avoidance proof using asynchronous renewals}\label{sec:asynchronous}
The ordered-probe theorem in the main argument yields the logarithmic
adjacent-arm bound. We record a different proof of the weaker exponent
estimate. It extends many pairs of asynchronously chosen renewal prefixes
to a common height and uses the telescoping strip identity there. This
argument uses the earlier localization and seam estimates, but is
independent of the ordered-probe theorem and fast-path amplification.
It makes no assumption that the two original paths share a renewal level.



The obstacle lemma below uses high moments of recoverable joining lines,
in contrast to the second-moment calculation giving the exact
\(H^{-1/4}\) bound in Lemma~\ref{lem:08-sharp-corridor}; this supplies a distinct
sewing argument for the avoidance proof.

Let \(S_2(H)\) be the probability that two independent infinite renewal
walks, launched vertically from neighboring horizontal bottom ports,
are disjoint through their respective first visits to level \(H\).
This event does not require the walks to have a common renewal level.
Let \(D(j,j)\) denote the unnormalized mass of two disjoint height-\(j\)
bridges from these ports.  Joining their starts by a fixed one-band
exterior arch, and translating to fix the first terminal port, produces
an arch counted at height \(j+1\) but not at height \(j\).  The input
is recoverable with bounded multiplicity.  The strip winding identity
therefore gives
\begin{equation}
                  D(j,j)\le C(B_j-B_{j+1}).          \label{eq:09-diagonal}
\end{equation}

\Needspace{5\baselineskip}
\begin{lemma}[Sewing past an obstacle]\label{09_fast_paths:obstacle}
Fix \(D<\infty\) and \(\delta>0\).  A port with abscissa \(x\)
on level \(u\) is to be joined to level \(v\), where
\(0\le v-u\le H\).  Suppose the obstacle in this slab lies to the
left of \(x+DH\), and in its first \(\min(\delta H,v-u)\) bands
lies to the left of \(x-\delta H\).  The mass of bridges to a free
endpoint at level \(v\) avoiding that obstacle is at least
\(H^{-1/4-o(1)}\), uniformly over these data.  The reflected statement
also holds.
\end{lemma}
\begin{proof}
For \(v-u\le\delta H/2\), use a vertical tube and
Lemma~\ref{lem:08-free-tube}.  Zero gap uses the empty path, and bounded
positive gaps use fixed elementary bridges.  For larger gaps construct
a corridor that initially rises vertically, moves far right within a
small fraction of the first \(\delta H\) bands, then rises beyond
\(x+(D+1)H\) to the target.  It consists of allowed normal directions:
up, right-up at \(30^\circ\), alternating right-down and right-up at
\(-30^\circ,30^\circ\), and finally up.  Consecutive directions have
positive dot product.  Choose each zigzag segment a small fixed
multiple of \(\delta H\), with equal lengths in alternating pairs,
and a sufficiently large fixed number of pairs to pass the obstacle.
All lengths are bounded above and below by positive multiples of
\(H\); the constants may depend on \(D,\delta\).  The corridor has
strict slack from the forbidden boundaries except at its endpoint walls.

At each joint vary a cut perpendicular to the next direction over a
small fixed-proportion range.  Before a corner \(C\) insert a parallel
joint \(F=C-\ell n_1\), where \(n_1\) is the current unit normal
and \(\ell\) is a sufficiently small fixed multiple of \(H\).
The stages are long parallel-to-parallel tubes and short corner pieces
from the \(F\)-cut, normal to \(n_1\), to the \(C\)-cut, normal to
\(n_2\), where \(n_1\cdot n_2=1/2\).  Each carries mass
\(\ge cH^{-1/4}\).  For long stages this is
Lemma~\ref{lem:08-free-tube}.  For a corner use the convex lattice
domain between the start line and a parallel far line with
\(n_1\cdot(p-C)\) near \(\ell/100\), bounded also by
\[
       -(.5+.02)\ell\ <\ (n_1-n_2)\cdot(p-C)\ <\ .02\ell,
\]
with lattice rounding.  This domain contains a thin normal tube past
the target cut if starting errors, cut ranges and widths are small
relative to \(\ell\).  Convex cutting in
Lemma~\ref{lem:08-localization} pays the short piece by the mass of
that tube.

The corner piece stays within \(O(\ell)\) of \(C\); its exit satisfies
\(n_1\cdot(p-F)>.4\ell\).  The next long tube, from that actual
exit, has width small enough that throughout it
\(n_1\cdot(p-F)>.3\ell\).  Choose the successive short scales
increasing so that accumulated errors are small fractions of the next
scale, while the final scale is still tiny compared with the original
segments.  There are only finitely many segments, so these choices are
consistent, and every cut range still has \(\asymp H\) levels.
Adjacent segments are separated at their common cut: before switching
at \(F\), the incoming tube lies strictly on its near side by the
positive dot product.  Nonconsecutive segments are separated by the
small widths.  Thus concatenated paths are simple and avoid the obstacle.

We check the multiplicity of these movable joints.  Only the two
adjoining stages visit the possible port region of a joint.  At \(F\)
ports are much closer to \(F\) than \(\ell\), clear of the outgoing
long tube; at \(C\) they are within \(O(\ell)\), clear of the incoming
long tube.  Thus split times for different joints occur in their fixed
temporal order in every representation.  If two representations split
differently at one joint, the intervening arc is a bridge between the
two cut lines.  It also remains strictly between the neighboring cut
ranges: at \(F\) it lies in the cut band and the incoming restricted
tube, hence very close to \(F\); at \(C\) the outgoing restricted
tube gives \(n_1\cdot(p-F)>.3\ell\).

Fix a base representation and select \(p\) other admissible levels
at one joint.  Before and after the base level these subdivide the
adjoining stages into bridges, one per ordered gap.  The end pieces
have bounded mass, apart from polynomial placement, by convex winding
in their endpoint half-planes with distant walls.  For example, the
piece before the earliest selected split remains in its near half-plane;
any change of split time at the preceding joint also lies strictly in
that half-plane by the preceding separation argument.  The same holds
after the last split.  Lemma~\ref{lem:08-seam-moments} therefore gives
\(C_pH^{C+3p/4}\) for the joint's possible-level count, with \(C\)
independent of \(p\).  A fixed level has a unique port and split time,
by the local single-crossing constraints.  High fixed moments discard
counts exceeding \(H^{3/4+\zeta}\) with arbitrary polynomially small
mass, for every \(\zeta>0\).

For \(m\) stages the construction sum is at least
\(cH^{m-1-m/4}\).  There are \(m-1\) joints, so division by their
multiplicities leaves
\(cH^{-1/4-(m-1)\zeta}\).  The discarded part is negligible relative
to the construction sum by choosing sufficiently high moments.  Since
\(m\) is fixed and \(\zeta\) arbitrary, this proves the lemma.
\end{proof}

\subsection{Horizontal displacement of an irreducible}
For one random irreducible bridge let \(J\) be its height and \(W\)
its largest absolute horizontal displacement from its starting port.
The renewal law already gives \(\Prob(J>x)\le C(1+x)^{-3/4}\).

\begin{lemma}[An irreducible displacement bound]\label{09_fast_paths:jump-tail}
For \(x\ge1\),
\[
 \Prob(W>x)\le Cx^{-3/4},\qquad
 \E\min\{1,(J+W)/x\}\le Cx^{-3/4}.
\]
\end{lemma}
\begin{proof}
Choose a large fixed \(A\), to be specified below. We prove the bound
for \(x\ge2A\); increasing the final constant covers \(1\le x<2A\).
Put \(r=\lfloor x/A\rfloor\ge1\), and write
\(p=\Prob(W>x,J\le r)\).  Let \(Z\) count renewal levels \(u\le r\) whose next irreducible
has \(W>x\) and \(J\le r\). Independence of the next increment gives
\[
 \E Z=p\sum_{u=0}^r B_u\asymp r^{3/4}p.
\]
For an ordered pair of distinct marks, expose the first marked jump
through its completion. Its terminal height is at most \(2r\).
Conditional on this history, the expected number of later marks born
at height at most \(r\) is zero if the endpoint exceeds \(r\), and
otherwise is at most \(p\sum_{v=0}^r B_v\le Cr^{3/4}p\).
Consequently
\[
 \E[Z(Z-1)]\le 2Cr^{3/4}p\,\E Z,
 \qquad \E Z^2\le Cr^{3/4}p(1+r^{3/4}p).
\]
Cauchy--Schwarz therefore gives
\(\Prob(Z>0)\ge c\min(1,r^{3/4}p)\).

On \(Z>0\), stop immediately after completing the first marked jump.
This is a stopping index in the irreducible sequence, at height at most
\(2r\); the fresh continuation is independent of the exposed history.
It forces an excursion farther than \(x/2\) from the original root:
either its start is already that far away, or its displacement forces
this during the jump.  After detection the expected number of further
renewals within another \(r\) bands is \(\asymp r^{3/4}\).
Every resulting bridge, ending by \(3r\), has that excursion.  Summing
the lateral loss in Lemma~\ref{lem:08-localization} gives total mass
at most \(Cr^2x^{-5/4}\).  Dividing by the continuation count bounds
the detection probability by \(C(r/x)^{5/4}\).  Take \(A\) large
enough to exclude the value 1 in the preceding minimum.  It follows
that \(p\le C_Ax^{-3/4}\); adding \(\Prob(J>r)\) proves the
first claim.  The union bound for \(J+W\) and integration of its tail
give the second:
\[
 \E\min(1,(J+W)/x)
 =x^{-1}\int_0^x\Prob(J+W>u)\,du\le Cx^{-3/4}.
\]
\end{proof}


\begin{theorem}[Adjacent-arm avoidance]\label{thm:asynchronous-avoidance}
For every \(\varepsilon>0\), \(S_2(H)\le C_\varepsilon
H^{-3/4+\varepsilon}\).  Equivalently,
\begin{equation}
                       S_2(H)\le H^{-3/4+o(1)}.
                 \label{eq:asynchronous-avoidance}
\end{equation}
\end{theorem}
\begin{proof}
Fix small \(\varepsilon>0\), put \(p_*=3/4-\varepsilon\), and
use strong induction with bound \(KH^{-p_*}\).  Choose a large fixed
\(L\), then set \(H_0=\lceil H/L\rceil\).  Stop each path at its
first renewal at or above \(H_0\).  On survival through \(H\), the
event that either prefix has sum of \(J+W\) greater than \(H/8\)
has probability at most
\begin{equation}
                      C_\varepsilon K L^{-\varepsilon}H^{-p_*}.
                                                     \label{eq:09-initial-cost}
\end{equation}
Indeed, group jumps starting below \(H_0\) into dyadic height bins
\([j,2j)\), with a separate bin for height zero.  For one path expose
the other path, and stop the first at its first renewal at or above
\(j\).  Survival through first visits to \(j\) is measurable then.
The expected subsequent sum of charges
\(\min(1,(J+W)/H)\) from starting heights in that bin is at most
\(Cj^{3/4}H^{-3/4}\), uniformly in the overshoot; if it skips the
bin the contribution is zero.  Multiply by \(S_2(j)\) and sum:
\[
 CH^{-3/4}\left[1+
       K\sum_{j\le H_0}^{\rm dyadic}j^\varepsilon\right]
                 \le C_\varepsilon K L^{-\varepsilon}H^{-p_*}.
\]
If a sum exceeds \(H/8\), its truncated charge sum is at least
\(1/8\), so this proves Equation~\eqref{eq:09-initial-cost}.
Choose \(L\) so this costs less than one quarter of the induction
budget.

Fix two stopped prefixes having these size bounds and surviving through
\(H_0\). Their future irreducibles remain independent. Under the law
of these two continuations, we shall exclude an event of arbitrarily
small prescribed probability, uniformly over the fixed prefixes.
Outside that event, continued survival through \(H\) supplies at least
\(cH^{3/4}\) suitable renewals on each side, in \([\eta H,3H/4]\),
such that every cross-pair can be extended to a disjoint equal-height
pair at cost \(H^{-1/4-o(1)}\).

We first retain many renewals in this interior height window while
controlling their horizontal range.
Take the next \(\lceil c'H^{3/4}\rceil\) jumps after either stop.
Lemma~\ref{09_fast_paths:jump-tail}, applied to their truncated heights,
shows that all these renewals lie within an additional \(H/2\) bands
except with probability \(O(c')\).  Renewals within \(\eta H\) of
the stop have expected number \(O((1+\eta H)^{3/4})\), and can be
discarded.  The full first-hit paths to \(H\) have horizontal range
at most \(D_0H\), apart from conditional probability
\(O(D_0^{-3/4})\): sum the truncated charges of jumps starting below
\(H\), whose expected number is \(O(H^{3/4})\).

For right renewals expose also the full left first-hit path.  Its
future is independent of the newly sampled right future.  At height
\(u\), let \(R(u)\) be its rightmost abscissa within
\([u,u+\delta H]\).  Lemma~\ref{lem:08-variation} implies
\[
 \sum_{j\ge\eta H}\|b_j\|_\infty
      \le C\sum_{j\ge\eta H}(B_j-B_{j+1})
      \le CB_{\lfloor\eta H\rfloor}.
\]
Thus the expected number of right renewals, at least \(\eta H\)
above their stop, with \(|x-R(u)|\le2eH\), is at most
\(C_\eta(e+H^{-1})H^{3/4}\).  Discard also renewals whose path
up to the next renewal more than \(\delta H\) above them travels
horizontally by more than \(eH\).  Restarting there and using the
truncated-charge bound gives expected discarded count at most
\(C(\delta/e)^{3/4}H^{3/4}\).

On survival, the disjoint first-hit arcs have fixed left-right order
across the slab.  At a given height no point of the left arc can lie
to the right of the right arc's rightmost point: that region is in
the right component of its complement.  At a retained right renewal
the right arc in the next \(\delta H\) bands has abscissa at most
\(x+eH\).  Hence \(R(u)\le x+eH\); exclusion of
\(|x-R(u)|\le2eH\) forces \(R(u)<x-2eH\).  The reflected argument
holds for left renewals.  Choose first \(c'\), then \(\eta\), then
\(e\), and then \(\delta<e\), sufficiently small, and finally
\(D_0\) large.  Markov's inequality and a union bound, applied to
the two tests without conditioning either on success of the other,
leave \(cH^{3/4}\) renewals per side with any prescribed fixed small
failure probability.  Lemma~\ref{09_fast_paths:obstacle} applies to
every cross-pair: the shorter bridge is extended above its renewal
level, so it cannot meet its own prefix, while the other prefix is
the obstacle.  This proves the assertion.

Choose the just-described failure probability less than
\(\tfrac14L^{-p_*}\), with a harmless adjustment for the ceiling in
\(H_0\).  Its total contribution is then at most one quarter of
\(KH^{-p_*}\), by the induction bound on \(S_2(H_0)\).

Finally, under the original product renewal law, let \(Z_H\) count pairs
of disjoint renewal prefixes in the stated height and horizontal ranges
for which the obstacle lemma's clearance conditions allow the prefix of
lesser height to be extended to the other's height. Equal heights use the
empty extension. Set \(Z_H=0\) unless the first-hit
paths survive through \(H\). These geometric conditions depend only on
the two finite prefixes, and their extension masses have a common lower
bound \(a_H=H^{-1/4-o(1)}\). A pair of prefixes occurs with exactly the
product of its bridge weights. Dropping the subsequent survival condition
and appending every allowed extension therefore gives
\[
 a_H\E Z_H\le\sum_{\Gamma}w(\Gamma)M(\Gamma),
\]
where \(\Gamma\) is an output pair of disjoint equal-height bridges,
\(w(\Gamma)\) is the product weight, and \(M(\Gamma)\) counts its
input representations. The common output height lies in
\([\eta H,H]\). Equation~\eqref{eq:09-diagonal} telescopes
to total output mass \(O_\eta(H^{-1/4})\).  For a fixed output, an
input is recovered by choosing which bridge was extended and a
single-crossing seam in that bridge; its unextended companion fixes
the common height. The two end pieces in the seam decomposition are
strip bridges. Dropping avoidance with the companion bounds its mass by
\(B_j\le1\), and summing the common height costs only \(O(H)\),
independently of the moment order. Lemma~\ref{lem:08-seam-moments} bounds the
multiplicity by \(H^{3/4+\zeta}\) outside arbitrarily small
polynomial mass.  Raw seam counts are \(O(H)\), so these exceptional
contributions remain negligible.  Dividing by the extension mass gives
\[
                          \E Z_H\le H^{3/4+o(1)}.
\]
On the retained survival event, \(Z_H\ge c^2H^{3/2}\).
Its probability is therefore \(H^{-3/4+o(1)}\), less than the
remaining half of the induction budget for sufficiently large \(H\).
Increase \(K\) to cover the finitely many smaller heights.  The
induction proves the theorem.
\end{proof}
